\chapter{Experiments and Sensitivity Analyses}
\label{ch:experiments}

We execute a battery of Monte Carlo and empirical experiments to validate the
theoretical claims of Chapters~\ref{ch:foundations} and~\ref{ch:methodology}.
All synthetic experiments use $N_{\text{sim}}=100$ paths of length $T=250$
steps under CRN (\S\ref{sec:crn}) unless stated otherwise.  Figures are
generated by \texttt{generate\_all\_figures.py} and displayed in both the
static Academic Logs page and the live simulator dashboard.

% ---------------------------------------------------------------
\section{Experiment 1: Stationary MAB Baseline}
\label{sec:exp1}
% ---------------------------------------------------------------

\paragraph{Setup.}
Single regime $r_t\sim\Normal(0.08,\,0.15^2)$.  All agents begin with a
uniform Beta$(1,1)$ prior.  This environment tests convergence to the Kelly
Oracle under ideal stationary conditions and validates Theorem~\ref{thm:kelly_slln}.

\paragraph{Results.}
Figure~\ref{fig:stationary_fan} shows that Thompson Sampling and EXP3
converge toward the Oracle path by $T\approx 50$, after which Fan bands
tighten.  The Risk-Constrained CPPI achieves lower median wealth
(Table~\ref{tab:stationary_stats}) because its leverage cap truncates
the upside in the absence of any crash: this is the cost of survival
in a benign environment.  Crucially, \emph{all} agents achieve zero ruin
probability in the stationary regime, confirming that the sticky-prior
problem is harmless when the environment does not change.

\begin{figure}[h]
\centering
\includegraphics[width=0.88\textwidth]{Stationary_wealth_fan.png}
\caption{Monte Carlo asymptotic trajectory fan --- Stationary environment.
Solid lines are median paths; shaded bands span the 10th--90th percentile.}
\label{fig:stationary_fan}
\end{figure}

\begin{table}[h]
\centering
\caption{Terminal distribution statistics --- Stationary environment
($T=250$, $N_{\text{sim}}=100$).}
\label{tab:stationary_stats}
\renewcommand{\arraystretch}{1.2}
\begin{tabular}{lccc}
\toprule
\textbf{Agent} & \textbf{Median Wealth} & \textbf{Max Wealth} & \textbf{Ruin Risk} \\
\midrule
Vol-HMM (Proposed)     & $3{,}019.7\times$ & $152{,}469\times$ & $0.0\%$ \\
EXP3 (Fractional)      & $7{,}944.9\times$ & $189{,}655\times$ & $0.0\%$ \\
Softmax                & $1{,}257.3\times$ & $2{,}111{,}300\times$ & $0.0\%$ \\
Thompson Sampling      & $1{,}012.0\times$ & $1{,}038{,}292\times$ & $0.0\%$ \\
Naive Bayes            &   $909.5\times$   & $1{,}139{,}905\times$ & $0.0\%$ \\
Risk-Constrained CPPI  &    $34.9\times$   &     $1{,}447\times$   & $0.0\%$ \\
\bottomrule
\end{tabular}
\end{table}

% ---------------------------------------------------------------
\section{Experiment 2: SLLN Convergence Verification}
\label{sec:exp2}
% ---------------------------------------------------------------

\paragraph{Setup.}
$N_{\text{sim}}=50$ paths of length $T=5{,}000$ under the Oracle agent.
The theoretical limit $\E[\log(1+f^*r)]$ is computed by numerical quadrature
over $\Normal(0.08,0.15^2)$.

\paragraph{Results.}
The empirical mean log-growth rate converges to within $10^{-3}$ of the
theoretical limit by $T=5{,}000$, providing a computational confirmation of
Theorem~\ref{thm:kelly_slln}.  This simultaneously validates the log-space
\texttt{WealthTracker} implementation: \texttt{numpy.log1p} accumulation
produces no numerical overflow or underflow errors across the extended horizon.

% ---------------------------------------------------------------
\section{Experiment 3: Horizon Trade-offs and Fractional Kelly}
\label{sec:exp3}
% ---------------------------------------------------------------

\paragraph{Setup.}
Kelly multipliers $c\in[0.1, 2.0]$ are swept over $N_{\text{sim}}=1{,}000$
short-horizon ($T=50$) paths.  Two statistics are recorded: (i) median
terminal wealth, and (ii) probability of reaching a $+25\%$ target.

\paragraph{Results.}
Full Kelly ($c=1.0$) maximizes median terminal wealth, consistent with
Theorem~\ref{thm:kelly_slln}.  Half-Kelly ($c=0.5$) maximizes the probability
of reaching the $+25\%$ target, consistent with Thorp (1969).  Over-betting
($c>1.0$) strictly decreases both metrics beyond a threshold, confirming that
$f > f^*$ is dominated by $f^*$ in the long run.

% ---------------------------------------------------------------
\section{Experiment 4: Adversarial Shock --- Regime Break at $T/2$}
\label{sec:exp4}
% ---------------------------------------------------------------

\paragraph{Setup.}
Deterministic bull-to-bear structural break at $T_{\text{shock}}=125$.
Bull: $r_t\sim\Normal(0.08,\,0.15^2)$.
Bear: $r_t\sim t_3(-0.10,\,0.30^2)$ (Student-$t$ with 3 degrees of freedom).
This is the primary experiment motivating the proposed methods.

\paragraph{Results.}
Figure~\ref{fig:adversarial_fan} shows that Thompson Sampling, Naive Bayes,
and Softmax fail to deleverage rapidly after the structural break, resulting
in median terminal wealth well below 1.0 (principal loss).  EXP3 responds
more quickly but still incurs 43\% ruin risk (Table~\ref{tab:adversarial_stats}).

\begin{figure}[h]
\centering
\includegraphics[width=0.88\textwidth]{Adversarial_wealth_fan.png}
\caption{Monte Carlo trajectory fan --- Adversarial Shock at $T=125$.
Vol-HMM (green) and Risk-Constrained CPPI (magenta) diverge upward from
the Bayesian baseline agents immediately after the structural break.}
\label{fig:adversarial_fan}
\end{figure}

The Vol-HMM detects the regime switch within approximately $1.9$ steps
post-break (Table~\ref{tab:aii_sensitivity}), cuts leverage to near zero,
and achieves median terminal wealth $42.3\times$ with only $3\%$ ruin risk.
The Risk-Constrained CPPI achieves $0\%$ ruin risk by
Proposition~\ref{prop:cppi_bound}, at the cost of lower median wealth
($1.07\times$).

\begin{table}[h]
\centering
\caption{Terminal distribution statistics --- Adversarial Shock
($T=250$, $N_{\text{sim}}=100$).}
\label{tab:adversarial_stats}
\renewcommand{\arraystretch}{1.2}
\begin{tabular}{lccc}
\toprule
\textbf{Agent} & \textbf{Median Wealth} & \textbf{Max Wealth} & \textbf{Ruin Risk} \\
\midrule
Vol-HMM (Proposed)     & $42.28\times$      & $623.3\times$   & $3.0\%$  \\
Risk-Constrained CPPI  &  $1.07\times$      &  $22.9\times$   & $0.0\%$  \\
Softmax                &  $0.84\times$      & $178.0\times$   & $12.0\%$ \\
Naive Bayes            &  $0.73\times$      & $118.1\times$   & $12.0\%$ \\
Thompson Sampling      &  $0.78\times$      & $117.5\times$   & $16.0\%$ \\
EXP3 (Fractional)      & $0.000095\times$   &   $7.6\times$   & $43.0\%$ \\
\bottomrule
\end{tabular}
\end{table}

% ---------------------------------------------------------------
\section{Experiment 5: Slow Nonstationary Drift}
\label{sec:exp5}
% ---------------------------------------------------------------

\paragraph{Setup.}
Mean drifts linearly from $\mu_0=0.08$ to $\mu_1=-0.05$ over $T=250$ steps;
$\sigma=0.15$ constant.  No abrupt change-point exists; agents must unlearn
their prior gradually.

\paragraph{Results.}
Figure~\ref{fig:drift_fan} shows all agents track the drift during the early
positive-mean phase.  Thompson Sampling and Naive Bayes accumulate regret as
the mean crosses zero (approximately $T=144$) because Beta posteriors retain
positive-expectation mass from prior observations.  The Vol-HMM responds when
its CUSUM trigger fires as $\sigma_{\text{roll}}$ climbs beyond the bull
emission scale, cutting leverage before the mean turns fully negative.
Figure~\ref{fig:drift_regret} confirms that the Vol-HMM and CPPI accumulate
the least Oracle regret across the full horizon.

\begin{figure}[h]
\centering
\includegraphics[width=0.88\textwidth]{Drift_wealth_fan.png}
\caption{Monte Carlo trajectory fan --- Slow Gaussian Drift.}
\label{fig:drift_fan}
\end{figure}

\begin{figure}[h]
\centering
\includegraphics[width=0.88\textwidth]{Drift_regret_dynamics.png}
\caption{Cumulative Oracle regret --- Slow Drift.  Lower paths indicate
superior convergence to the growth-optimal benchmark.}
\label{fig:drift_regret}
\end{figure}

% ---------------------------------------------------------------
\section{Experiment 6: Kelly-Ruin Paradox in Heavy Tails}
\label{sec:kelly_ruin}
% ---------------------------------------------------------------

\paragraph{Setup.}
The Gaussian Kelly fraction $f^*=\mu/\sigma^2$ is computed from sample
moments, then applied to both a Gaussian and a Student-$t_3$ environment
with identical $\mu=0.08$ and $\sigma=0.15$.  $N_{\text{sim}}=1{,}000$
paths, $T=500$ steps.

\paragraph{Results.}
Ruin probability under the Gaussian: $\approx 0\%$.  Ruin probability under
the Student-$t_3$: $\approx 12\%$.  This confirms the Kelly-Ruin Paradox:
applying a formula derived under Gaussian assumptions to a heavy-tailed
environment causes ruin at a non-trivial rate, directly motivating the
tail-agnostic CPPI safety layer.

% ---------------------------------------------------------------
\section{Sensitivity Analyses}
\label{sec:sensitivity}
% ---------------------------------------------------------------

\subsection{Vol-HMM Transition Persistence $A_{ii}$}

\begin{table}[h]
\centering
\caption{Effect of HMM persistence $A_{ii}$ on detection lag and ruin
probability in the Adversarial Shock environment.}
\label{tab:aii_sensitivity}
\renewcommand{\arraystretch}{1.2}
\begin{tabular}{ccc}
\toprule
\textbf{Persistence $A_{ii}$} & \textbf{Mean Detection Lag (steps)} & \textbf{Ruin Risk} \\
\midrule
$0.99$ & $14.3$ & $32.4\%$ \\
$0.98$ & $8.7$  & $15.1\%$ \\
$\mathbf{0.95}$ \textbf{(selected)} & $\mathbf{1.9}$ & $\mathbf{0.0\%}$ \\
$0.80$ & $1.2$  & $0.0\%$ (excessive whipsawing) \\
\bottomrule
\end{tabular}
\end{table}

$A_{ii}=0.95$ is selected as the operating point.  Lower values
($A_{ii}=0.80$) achieve faster detection but cause excessive
whipsawing---frequent false bear-regime detections in the bull phase that
reduce gross returns unnecessarily.

\subsection{CUSUM Drift Parameter $k_{\mathrm{drift}}$}

\begin{table}[h]
\centering
\caption{Effect of CUSUM drift parameter on CAGR and false-positive
allocation drawdown.}
\label{tab:cusum_sensitivity}
\renewcommand{\arraystretch}{1.2}
\begin{tabular}{ccc}
\toprule
\textbf{$k_{\mathrm{drift}}$} & \textbf{CAGR (Survival Path)} & \textbf{False Positive Drawdown} \\
\midrule
$1.0$                                 & $4.5\%$  & High  \\
$\mathbf{0.5}$ \textbf{(selected)}    & $8.2\%$  & Low   \\
$0.2$                                 & Ruin-adjacent & Minimal \\
\bottomrule
\end{tabular}
\end{table}

\subsection{CPPI Multiplier $m$ and Floor $D_{\max}$}

Setting $m=3.0$ and $D_{\max}=20\%$ achieves the operating point used
throughout the main experiments: zero floor breaches on the adversarial
shock with median terminal wealth $1.07\times$.  Increasing $m$ raises
expected growth but also the probability of discrete-time floor breaches
during large single-period returns, consistent with the discrete-time
caveat in Proposition~\ref{prop:cppi_bound}.

% ---------------------------------------------------------------
\section{Experiment 7: Statistical Significance}
\label{sec:exp7}
% ---------------------------------------------------------------

\paragraph{Fisher's Exact Test.}
We test whether the Vol-HMM achieves a strictly lower ruin probability than
EXP3 under the Adversarial Shock environment.  The $2\times 2$ contingency
table (ruined/survived $\times$ Vol-HMM/EXP3) is tested via Fisher's Exact
Test (two-sided, \texttt{scipy.stats.fisher\_exact}).  With
$N_{\text{sim}}=100$, $p < 0.001$ consistently, confirming statistically
significant ruin-probability reduction.

\paragraph{Mann-Whitney U Test.}
Testing UCB vs.\ Naive Bayes terminal wealth in the stationary environment
with $N_{\text{sim}}=300$ paths: Mann-Whitney $U$ yields $p < 0.05$,
confirming statistically significant wealth outperformance under the
non-parametric test appropriate for log-normal distributions.  The
live dashboard reports Fisher's Exact $p$-values in real time on the
Math Report tab.

% ---------------------------------------------------------------
\section{Historical Empirical Validation (S\&P 500)}
\label{sec:historical}
% ---------------------------------------------------------------

All agents are evaluated on single-path historical log-returns for the
S\&P~500 across four distinct epochs, downloaded via \texttt{yfinance} and
cached as NumPy \texttt{.npy} arrays.  The CRN framework degenerates to
$N_{\text{sim}}=1$ for historical replay; the same agent \texttt{act()}
and \texttt{update()} loops are used to maintain ``one framework'' integrity.

\begin{table}[h]
\centering
\caption{Empirical agent performance --- 2008 Global Financial Crisis
(100 bps friction).}
\label{tab:historical_2008}
\renewcommand{\arraystretch}{1.2}
\begin{tabular}{lccc}
\toprule
\textbf{Agent} & \textbf{Terminal Wealth} & \textbf{Sharpe} & \textbf{Sortino} \\
\midrule
Buy-and-Hold (benchmark) & $0.62\times$ & $-0.45$ & $-0.32$ \\
Vol-HMM (Proposed)       & $0.88\times$ & $+0.12$ & $+0.45$ \\
Risk-Constrained CPPI    & $0.80\times$ & $+0.05$ & $+0.82$ \\
Thompson Sampling        & $0.45\times$ & $-0.85$ & $-0.60$ \\
\bottomrule
\end{tabular}
\end{table}

The CPPI achieves the highest Sortino ratio by explicitly truncating the left
tail (Proposition~\ref{prop:cppi_bound}).  The Vol-HMM detects the volatility
regime shift of September 2008 and March 2020 within approximately 2 trading
days, significantly outperforming the unconstrained Bayesian baselines.

% ---------------------------------------------------------------
\section{Transaction Fee Audit}
\label{sec:fee_audit}
% ---------------------------------------------------------------

High-turnover agents (EXP3, UCB) rebalance aggressively at every step.
Under 25--50 bps friction, their capital erosion reaches up to $15\%$ per
annum (annualized from $T=250$ paths with $252$ trading days).  The Vol-HMM
maintains a stable regime belief between detected transitions, producing
lower turnover and significantly lower cumulative fees.  The dashboard
``Math Report'' tab reports both gross terminal wealth and net-of-fees
terminal wealth, making this trade-off directly observable during the defense.

\begin{figure}[h]
\centering
\includegraphics[width=0.88\textwidth]{Adversarial_regret_dynamics.png}
\caption{Cumulative Oracle regret --- Adversarial Shock.  The sharp upward
inflection at $T=125$ for the Bayesian baselines marks the onset of the
sticky-prior lag following the structural break.}
\label{fig:adversarial_regret}
\end{figure}
